% problem_id: S001-proposition-blowup-split
% source_id: S001 (Stacks Project)
% source_file: derham.tex
% statement_locator: derham.tex lines 2970--2975
% proof_locator: derham.tex lines 2977--3094
% env: proposition
% pairing: tight (verified)
% status: complete (statement + author proof verbatim + reason + locators)
%
\begin{proposition}
\label{proposition-blowup-split}
With notation as in More on Morphisms, Lemma \ref{more-morphisms-lemma-blowup}
the map $\Omega^\bullet_{X/S} \to Rb_*\Omega^\bullet_{X'/S}$
has a splitting in $D(X, (X \to S)^{-1}\mathcal{O}_S)$.
\end{proposition}

\begin{proof}
Consider the triangle constructed in
Lemma \ref{lemma-distinguished-triangle-blowup}.
We claim that the map
$$
Rb_*(\Omega^\bullet_{X'/S}) \oplus i_*\Omega^\bullet_{Z/S} \to
i_*Rp_*(\Omega^\bullet_{E/S})
$$
has a splitting whose image contains the summand $i_*\Omega^\bullet_{Z/S}$.
By Derived Categories, Lemma \ref{derived-lemma-split} this will show that
the first arrow of the triangle has a splitting which vanishes on
the summand $i_*\Omega^\bullet_{Z/S}$ which proves the lemma.
We will prove the claim by decomposing $Rp_*\Omega^\bullet_{E/S}$
into a direct sum where the first piece corresponds to
$\Omega^\bullet_{Z/S}$ and the second piece can be lifted
through $Rb_*\Omega^\bullet_{X'/S}$.

\medskip\noindent
Proof of the claim. We may decompose $X$ into open and closed subschemes
having fixed relative dimension to $S$, see
Morphisms, Lemma \ref{morphisms-lemma-smooth-omega-finite-locally-free}.
Since the derived category $D(X, f^{-1}\mathcal{O})_S)$ correspondingly
decomposes as a product of categories, we may assume $X$ has
fixed relative dimension $N$ over $S$. We may decompose
$Z = \coprod Z_m$ into open and closed subschemes of relative
dimension $m \geq 0$ over $S$. The restriction $i_m : Z_m \to X$ of
$i$ to $Z_m$ is a regular immersion of codimension $N - m$, see Divisors, Lemma
\ref{divisors-lemma-immersion-smooth-into-smooth-regular-immersion}.
Let $E = \coprod E_m$ be the corresponding decomposition, i.e.,
we set $E_m = p^{-1}(Z_m)$. If $p_m : E_m \to Z_m$ denotes the
restriction of $p$ to $E_m$, then we have a canonical isomorphism
$$
\tilde \xi_m :
\bigoplus\nolimits_{t = 0, \ldots, N - m - 1}
\Omega^\bullet_{Z_m/S}[-2t]
\longrightarrow
Rp_{m, *}\Omega^\bullet_{E_m/S}
$$
in $D(Z_m, (Z_m \to S)^{-1}\mathcal{O}_S)$
where in degree $0$ we have the canonical map
$\Omega^\bullet_{Z_m/S} \to Rp_{m, *}\Omega^\bullet_{E_m/S}$.
See Remark \ref{remark-projective-space-bundle-formula}.
Thus we have an isomorphism
$$
\tilde \xi :
\bigoplus\nolimits_m
\bigoplus\nolimits_{t = 0, \ldots, N - m - 1}
\Omega^\bullet_{Z_m/S}[-2t]
\longrightarrow
Rp_*(\Omega^\bullet_{E/S})
$$
in $D(Z, (Z \to S)^{-1}\mathcal{O}_S)$
whose restriction to the summand
$\Omega^\bullet_{Z/S} = \bigoplus \Omega^\bullet_{Z_m/S}$ of the source
is the canonical map $\Omega^\bullet_{Z/S} \to Rp_*(\Omega^\bullet_{E/S})$.
Consider the subcomplexes $M_m$ and $K_m$ of the complex
$\bigoplus\nolimits_{t = 0, \ldots, N - m - 1} \Omega^\bullet_{Z_m/S}[-2t]$
introduced in Remark \ref{remark-projective-space-bundle-formula}.
We set
$$
M = \bigoplus M_m
\quad\text{and}\quad
K = \bigoplus K_m
$$
We have $M = K[-2]$ and by construction the map
$$
c_{E/Z} \oplus \tilde \xi|_M :
\Omega^\bullet_{Z/S} \oplus M
\longrightarrow
Rp_*(\Omega^\bullet_{E/S})
$$
is an isomorphism (see remark referenced above).

\medskip\noindent
Consider the map
$$
\delta : \Omega^\bullet_{E/S}[-2] \longrightarrow \Omega^\bullet_{X'/S}
$$
in $D(X', (X' \to S)^{-1}\mathcal{O}_S)$ of
Lemma \ref{lemma-gysin-via-log-complex}
with the property that the composition
$$
\Omega^\bullet_{E/S}[-2] \longrightarrow \Omega^\bullet_{X'/S}
\longrightarrow
\Omega^\bullet_{E/S}
$$
is the map $\theta'$ of Remark \ref{remark-cup-product-as-a-map} for
$c_1^{dR}(\mathcal{O}_{X'}(-E))|_E) = c_1^{dR}(\mathcal{O}_E(1))$.
The final assertion of Remark \ref{remark-projective-space-bundle-formula}
tells us that the diagram
$$
\xymatrix{
K[-2] \ar[d]_{(\tilde \xi|_K)[-2]} \ar[r]_{\text{id}} &
M \ar[d]^{\tilde x|_M} \\
Rp_*\Omega^\bullet_{E/S}[-2] \ar[r]^-{Rp_*\theta'} &
Rp_*\Omega^\bullet_{E/S}
}
$$
commutes. Thus we see that we can obtain the desired splitting of
the claim as the map
\begin{align*}
Rp_*(\Omega^\bullet_{E/S})
& \xrightarrow{(c_{E/Z} \oplus \tilde \xi|_M)^{-1}}
\Omega^\bullet_{Z/S} \oplus M \\
& \xrightarrow{\text{id} \oplus \text{id}^{-1}}
\Omega^\bullet_{Z/S} \oplus K[-2] \\
& \xrightarrow{\text{id} \oplus (\tilde \xi|_K)[-2]}
\Omega^\bullet_{Z/S} \oplus Rp_*\Omega^\bullet_{E/S}[-2] \\
& \xrightarrow{\text{id} \oplus Rb_*\delta}
\Omega^\bullet_{Z/S} \oplus Rb_*\Omega^\bullet_{X'/S}
\end{align*}
The relationship between $\theta'$ and $\delta$ stated above
together with the commutative diagram involving $\theta'$, $\tilde \xi|_K$,
and $\tilde \xi|_M$ above are exactly what's needed to
show that this is a section to the canonical map
$\Omega^\bullet_{Z/S} \oplus Rb_*(\Omega^\bullet_{X'/S}) \to
Rp_*(\Omega^\bullet_{E/S})$ and the proof of the claim is complete.
\end{proof}
